% Replace the existing purity bullet point in the XGBoost Data Generation section with this expanded version.
% This goes inside the \begin{itemize} block, replacing the original \item \textbf{Tracklet purity threshold} bullet.

    \item \textbf{Tracklet purity threshold}. Each tracklet is assigned a ground truth identity based on the majority identity among its detections. The purity of a tracklet is defined as the fraction of detections that belong to this majority identity. For example, a tracklet with 100 frames where 70 belong to player \#5 and 30 belong to player \#12 has a purity of 0.70, and is assigned the identity of player \#5. Two tracklets are labeled as a positive pair (should merge) when both have an assigned identity and these identities match. Tracklets with a purity below the threshold receive no identity assignment, and any pair involving such a tracklet is labeled as negative regardless of the other tracklet's identity.

    The purity threshold introduces a tradeoff between label quality and distribution matching. A strict threshold (e.g., 1.0) produces clean labels, since only tracklets that belong entirely to one player contribute to positive pairs. However, it also excludes many tracklets from the training data and creates a distribution mismatch with inference: the splitter misses approximately 85\% of identity switches (Section~\ref{sec:splitter_eval}), so the connector receives impure tracklets at inference time. These impure tracklets contain blended features: their mean ReID embedding, jersey predictions, and team predictions reflect a mixture of two players. A model trained exclusively on pure tracklets has never encountered these blended feature patterns and may not generalize to them.

    A permissive threshold (e.g., 0.6) accepts noisier tracklets, producing more positive pairs whose features better match the impure tracklets seen at inference time. The cost is some label noise: a tracklet that is 60\% player \#5 and 40\% player \#12 gets labeled as player \#5, so a pair with another player \#5 tracklet receives a positive label even though the features are partially contaminated by a second player. We test purity thresholds of 0.6, 0.8, and 1.0.
